\section{Vi điều khiển ESP32}

\subsection{Tổng quan ESP32}

ESP32 là dòng vi điều khiển SoC (System on Chip) do Espressif Systems phát triển,
được thiết kế cho các ứng dụng IoT và thiết bị nhúng. Chip tích hợp bộ xử lý
Xtensa LX6 dual-core chạy ở tần số lên đến 240 MHz, cùng với WiFi 802.11 b/g/n
và Bluetooth 4.2/BLE trên một chip duy nhất. Với giá thành thấp (khoảng 2--5 USD)
và hệ sinh thái phong phú, ESP32 trở thành lựa chọn phổ biến cho các dự án
nhúng từ hobby đến thương mại.

Dự án sử dụng board ESP32-CAM --- một biến thể tích hợp sẵn khe cắm camera
và hỗ trợ PSRAM 4 MB mở rộng, cho phép xử lý ảnh và buffer dữ liệu lớn
mà không bị giới hạn bởi 520 KB SRAM nội bộ.

\begin{table}[htbp]
\centering
\caption{Thông số phần cứng ESP32-CAM}
\label{tab:esp32-hardware}
\begin{tabular}{>{\raggedright\arraybackslash}p{0.3\linewidth} >{\raggedright\arraybackslash}p{0.6\linewidth}}
\toprule
\rowcolor{headerblue} \textcolor{white}{\textbf{Thành phần}} & \textcolor{white}{\textbf{Thông số}} \\
\midrule
\rowcolor{rowgray} Board & ESP32-CAM với PSRAM \\
Vi xử lý & Xtensa LX6 dual-core, 240 MHz \\
\rowcolor{rowgray} Camera & OV5640, JPEG, QVGA (320$\times$240) mặc định, hỗ trợ đến UXGA (1600$\times$1200) \\
Microphone & INMP441 (I2S digital), 16 kHz, mono \\
\rowcolor{rowgray} Kết nối & WiFi 802.11 b/g/n, 2.4 GHz \\
Bộ nhớ & 520 KB SRAM + 4 MB PSRAM \\
\rowcolor{rowgray} Flash & 4 MB SPI Flash \\
GPIO & 16 chân GPIO đa chức năng \\
\rowcolor{rowgray} Nguồn & 5V qua USB hoặc pin, 3.3V logic \\
\bottomrule
\end{tabular}
\end{table}

\begin{figure}[htbp]
\centering
\fbox{\parbox{0.85\linewidth}{\centering\vspace{2.5cm}
\textit{Thêm hình: Ảnh thực tế board ESP32-CAM --- mặt trước hiển thị chip ESP32,
khe cắm camera, ăng-ten WiFi. Mặt sau hiển thị khe cắm thẻ microSD
và PSRAM chip.}
\vspace{2.5cm}}}
\caption{Board ESP32-CAM}
\label{fig:esp32-cam-board}
\end{figure}

\subsection{Kiến trúc dual-core và FreeRTOS}

ESP32 sử dụng hai nhân xử lý Xtensa LX6 (gọi là Core 0 và Core 1),
cho phép chạy song song hai tác vụ độc lập. Hệ điều hành thời gian thực
FreeRTOS được tích hợp sẵn trong ESP-IDF (Espressif IoT Development Framework),
cung cấp cơ chế quản lý task, semaphore, queue và timer. Mỗi task được gán
vào một core cụ thể thông qua hàm \texttt{xTaskCreatePinnedToCore()}, đảm bảo
tác vụ thời gian thực (như thu âm) không bị preempt bởi các tác vụ khác.

\subsection{Camera OV5640}

OV5640 là cảm biến hình ảnh CMOS 5 megapixel của OmniVision, hỗ trợ nhiều
độ phân giải từ QQVGA (160$\times$120) đến QSXGA (2592$\times$1944).
Trong dự án, camera được cấu hình ở chế độ JPEG để giảm tải cho ESP32 ---
thay vì truyền raw pixel, camera nén ảnh trước khi gửi qua bus dữ liệu 8-bit
song song. Camera giao tiếp với ESP32 qua hai giao thức: bus dữ liệu song song
(8 chân D0--D7) cho truyền pixel, và I2C (SDA/SCL) cho cấu hình các thanh ghi
điều khiển như độ sáng, contrast, white balance.

\subsection{Microphone INMP441}

INMP441 là microphone MEMS (Micro-Electro-Mechanical Systems) digital của
InvenSense (TDK), tích hợp sẵn bộ chuyển đổi ADC và giao tiếp I2S. So với
microphone analog truyền thống cần thêm mạch khuếch đại và ADC bên ngoài,
INMP441 đơn giản hóa đáng kể thiết kế phần cứng. Microphone có độ nhạy
$-26$ dBFS, SNR 61 dBA, và dải tần đáp ứng 60 Hz -- 15 kHz, phù hợp cho
thu giọng nói trong môi trường bán hàng.

\begin{figure}[htbp]
\centering
\fbox{\parbox{0.85\linewidth}{\centering\vspace{2cm}
\textit{Thêm hình: Ảnh module microphone INMP441 --- hiển thị chip MEMS,
các chân kết nối (VDD, GND, SD, SCK, WS, L/R), kích thước nhỏ gọn.}
\vspace{2cm}}}
\caption{Module microphone INMP441}
\label{fig:inmp441-module}
\end{figure}

\subsection{Giao thức I2S}

I2S (Inter-IC Sound) là giao thức truyền thông nối tiếp đồng bộ được thiết kế
chuyên biệt cho dữ liệu âm thanh số, ban đầu được Philips (nay là NXP) đề xuất
vào năm 1986. Giao thức sử dụng 3 đường tín hiệu chính:

\begin{itemize}
\item \textbf{SCK (Serial Clock):} Xung nhịp đồng bộ dữ liệu, tần số bằng
  tần số lấy mẫu $\times$ số bit $\times$ số kênh. Với cấu hình 16 kHz,
  32-bit, stereo: SCK = 1.024 MHz.
\item \textbf{WS (Word Select):} Chọn kênh trái (WS = 0) hoặc kênh phải (WS = 1).
  Tần số WS bằng tần số lấy mẫu (16 kHz).
\item \textbf{SD (Serial Data):} Đường dữ liệu nối tiếp, truyền MSB first.
\end{itemize}

Microphone INMP441 hỗ trợ I2S với độ phân giải 18-bit thực tế, đóng gói trong
word 32-bit với dữ liệu nằm ở các bit cao. ESP32 hoạt động ở chế độ I2S Master
RX, cung cấp xung nhịp SCK và WS, nhận dữ liệu từ microphone qua SD.

\begin{figure}[htbp]
\centering
\fbox{\parbox{0.85\linewidth}{\centering\vspace{2cm}
\textit{Thêm hình: Giản đồ thời gian (timing diagram) giao thức I2S --- hiển thị
3 tín hiệu SCK, WS, SD với dữ liệu 32-bit mỗi kênh. Chú thích MSB first,
vùng dữ liệu 18-bit hữu ích và 14-bit padding.}
\vspace{2cm}}}
\caption{Giản đồ thời gian giao thức I2S}
\label{fig:i2s-timing}
\end{figure}

\subsection{Giao thức TCP}

TCP (Transmission Control Protocol) là giao thức tầng vận chuyển đảm bảo truyền
dữ liệu tin cậy, có thứ tự và không trùng lặp giữa hai thiết bị qua mạng IP.
TCP sử dụng cơ chế three-way handshake để thiết lập kết nối, sequence number
để đảm bảo thứ tự, và acknowledgment kết hợp retransmission để xử lý mất gói.

Dự án chọn TCP thay vì UDP cho cả stream audio (port 3002) và video (port 3001)
vì hai lý do chính: thứ nhất, mất mẫu âm thanh sẽ gây artifact trong quá trình
khử nhiễu DTLN vì mô hình yêu cầu đầu vào liên tục; thứ hai, mất frame JPEG
không thể khôi phục được do mỗi frame được nén độc lập. Tuy TCP có overhead
cao hơn UDP, nhưng trong mạng LAN WiFi với latency thấp, sự khác biệt này
không đáng kể.

\subsection{Giao thức WebSocket}

WebSocket là giao thức truyền thông full-duplex trên một kết nối TCP duy nhất,
được chuẩn hóa trong RFC 6455. Khác với HTTP truyền thống (request-response),
WebSocket cho phép server chủ động đẩy dữ liệu đến client mà không cần client
gửi request --- đặc tính quan trọng cho việc streaming real-time video và audio
từ ESP32 đến trình duyệt.

Kết nối WebSocket bắt đầu bằng một HTTP Upgrade handshake, sau đó chuyển sang
giao thức WebSocket với overhead header chỉ 2--14 bytes (so với hàng trăm bytes
của HTTP header), giảm đáng kể bandwidth cho streaming liên tục.
